%% =====================================================================
%%  Avant-propos — édition 2027-2028
%% =====================================================================
\thispagestyle{otFront}
\phantomsection\pdfbookmark[0]{Avant-propos}{bm-avant-propos}
\otPageTitre{Avant-propos}{}

\begin{minipage}{128mm}
{\fontsize{14}{20}\selectfont\itshape\color{otNavy}
Ce cahier a été conçu pour t'accompagner tout au long de ton année de
Terminale, en spécialité mathématiques.\par}
\vspace{7mm}

\fontsize{11}{16}\selectfont
Cette \textbf{édition 2027-2028} suit le nouveau programme de Terminale, applicable à la
rentrée 2027, et la définition de l'épreuve du baccalauréat publiée en septembre 2026.
Elle réunit plus de 280 exercices, répartis en 23 chapitres et classés par niveau de
difficulté. Tous sont accompagnés d'un corrigé détaillé, regroupé en fin d'ouvrage.\par\vspace{3.5mm}

Chaque chapitre suit la même progression. Les exercices de
\textbf{Découverte} vérifient que les notions du cours sont en place ; ceux
d'\textbf{Entraînement} installent les méthodes incontournables ; les
\textbf{Problèmes}, enfin, te préparent aux sujets du baccalauréat.\par\vspace{3.5mm}

Pour terminer, dix \textbf{sujets type bac} complets te placent en conditions
d'examen : quatre exercices, quatre heures, vingt points dont deux pour la maîtrise
de la langue. Leurs corrigés détaillent le barème, question par question.\par\vspace{3.5mm}

Les corrigés ne se contentent pas de donner la réponse : ils indiquent, à
chaque étape, la méthode utilisée. Prends le temps de chercher avant de les
consulter : c'est en cherchant que l'on progresse.\par\vspace{3.5mm}

Les nouveautés du programme sont présentées page suivante et signalées dans le livre
par la mention \otNouveau. Le tableau de bord te permet de suivre ta progression
chapitre par chapitre.\par\vspace{9mm}

{\otMedium\fontsize{10}{13}\selectfont Bon travail, et bonne réussite au bac !}\par
\vspace{3mm}
{\otLight\fontsize{9}{12}\selectfont\color{otMuted}L'équipe pédagogique OptiTAB}\par
\vspace{8mm}
\includegraphics[height=9mm]{icone-couleur}
\end{minipage}
\clearpage
